\chapter{Transforms, ODEs, and state space}
\label{app:transforms}

The methods of this appendix apply to the convolution case $K(t,u)=\phi(t-u)$; their
failure off it is what makes excitation covariates hard (Ch.~\ref{ch:excitation}).

\section{Laplace, Fourier, and the convolution theorem}

The transforms turn convolution into multiplication,
$\Lap\{\phi*\xi\}=\Lap\{\phi\}\Lap\{\xi\}$. Applying $\Lap$ to $\xi=s+\phi*\xi$,
\begin{equation}
\Lap\{\xi\}=\frac{\Lap\{s\}}{1-\Lap\{\phi\}},\qquad
\hat\xi(\omega)=R(\omega)\hat s(\omega),\quad R=\frac{1}{1-\hat\phi}.
\end{equation}
$R$ is the \emph{transfer function} (the resolvent of Appendix~\ref{app:volterra} in
the frequency domain). For the exponential kernel $\hat\phi=\kappa\theta/(\theta+i\omega)$,
so $R(\omega)=(\theta+i\omega)/((1-\kappa)\theta+i\omega)$ --- a one-pole filter
whose pole is the ODE rate below.

\begin{intuitionbox}
Transforms re-express $\xi$ in pure oscillations; a shift-invariant system scales
each independently, so the tangled convolution becomes a per-frequency multiply. A
covariate that changes the gain over time breaks shift-invariance --- hence no single
transfer function in Ch.~\ref{ch:excitation}, and one must return to the time domain.
\end{intuitionbox}

\section{From exponential kernel to a first-order ODE}

Differentiating $y=\phi*\xi$ for $\phi=\kappa\theta e^{-\theta t}$ gives
\begin{equation}
y'(t)=(\kappa-1)\theta\,y(t)+\kappa\theta\,s(t),\qquad \xi=s+y,
\end{equation}
a stable first-order ODE forced by the baseline (pole $(\kappa-1)\theta<0$).

\begin{theorem}[realisation]\label{thm:appode-rational}
The convolution MBPP equation is equivalent to a finite linear constant-coefficient
ODE iff $\phi$ has a strictly proper rational Laplace transform --- equivalently,
$\phi$ is a finite exponential-polynomial.
\end{theorem}
A sum of $Q$ exponentials gives the $Q$-state realisation
$\dot u=Au+a\,s$, $\xi=s+\mathbf 1^\top u$ (Ch.~\ref{ch:multitimescale}); the
multivariate kernel gives the $M^2$-state system of Ch.~\ref{ch:solvers}.

\section{Time-varying systems and the fractional case}

For excitation covariates the realisation is the same in form but with a
time-varying gain (Ch.~\ref{ch:excitation}): a \emph{linear time-varying} ODE,
integrable step by step (exactly per-regime for piecewise-constant covariates) but
with no closed matrix exponential when the $A(t)$ do not commute --- the
differential-equations echo of ``no constant transfer function.'' For a power-law
kernel the transform is \emph{not} rational, no finite ODE exists, and the resolvent
is a Mittag--Leffler function $E_{\alpha,\alpha}$ (the dynamics are fractional);
$E_{1,1}=\exp$ recovers the exponential case. In practice one approximates the
power law by a sum of exponentials (restoring a finite ODE) or solves by Volterra
quadrature.
